\section{Equilibrio idrostatico}

\subsection{Equilibrio newtoniano}

% Frame 1
\begin{frame}[t]{Capitolo 4: mappa degli argomenti}
\small
Il capitolo collega la microfisica di materia densa alle configurazioni di equilibrio gravitazionale.
\begin{enumerate}
\item Dall'equilibrio idrostatico newtoniano alle TOV.
\item Argomento dimensionale di Landau: masse massime e nane bianche.
\item Invarianza di scala delle TOV e limiti causali.
\item Fermioni interagenti e stelle bosoniche.
\end{enumerate}
\medskip
\textit{Collegamenti con gli appunti precedenti}: capitoli 2 (RG e TOV), 3 (EOS di Fermi), 8 (materia strange e bag), 9 (stelle ibride e terza famiglia).
\medskip
Convenzioni: $c=\hbar=1$ salvo indicazione; $G=m_{\rm Pl}^{-2}$ con massa di Planck \emph{non ridotta}.
\end{frame}

% Frame 2
\begin{frame}[t]{Massa racchiusa in una sfera: equazioni (4.1)--(4.2)}
\small
Per una distribuzione sferica $\rho(r)$, la massa racchiusa nella sfera di raggio $r$ è
\[
m(r)=\int_{|\mathbf x|<r}\rho(\mathbf x)\,d^3x
=4\pi\int_0^r\rho(r')\,r'^2\,dr'.
\]
Il teorema fondamentale del calcolo fornisce
\[
\frac{dm}{dr}=4\pi r^2\rho(r),\qquad m(0)=0,\quad m(R)=M.
\]
$\rho$ è densità di massa in unità newtoniane; $m(r)$ è massa \emph{integrata}, non densità locale. La derivazione vale anche per $\rho$ non uniforme.
\end{frame}

% Frame 3
\begin{frame}[t]{Equilibrio idrostatico: da forze a equazioni (4.3)--(4.7)}
\small
Consideriamo un elemento di fluido di area $A$ e spessore $dr$:
\[
dm=A\rho(r)\,dr,\quad dF_G=-\frac{Gm(r)}{r^2}\,A\rho(r)\,dr.
\]
La differenza fra le forze di pressione sulle facce opposte è
\[
dF_P=A\bigl[P(r)-P(r+dr)\bigr]=-A\,dP.
\]
Ponendo $dF_P+dF_G=0$ (con versi opportuni):
\[
\frac{dP}{dr}=-\frac{Gm(r)\rho(r)}{r^2}.
\]
Condizioni tipiche: $P(0)=P_c$, $P(R)=0$ in un modello con superficie nel vuoto.
\end{frame}

% Frame 4
\begin{frame}[t]{Che cosa chiude il sistema di equilibrio?}
\small
Le due equazioni
\[
\frac{dP}{dr}=-\frac{Gm\rho}{r^2},\qquad
\frac{dm}{dr}=4\pi r^2\rho
\]
contengono le tre incognite $P(r)$, $m(r)$ e $\rho(r)$.
\medskip
Serve l'equazione di stato (EOS), ad esempio $P=P(\rho)$; fissando $P_c$ otteniamo una soluzione.
\medskip
La scelta dell'EOS è il punto in cui entrano i gradi di libertà microscopici (elettroni degeneri, neutroni, interazioni forti, quark, ecc.).
\end{frame}

\subsection{Equilibrio relativistico}

% Frame 5
\begin{frame}[t]{Fluido perfetto e conservazione in relatività generale}
\small
Adottiamo qui la metrica $(-,+,+,+)$ come nel libro e normalizzazione $u_\mu u^\mu=-1$.
\[
T^{\mu\nu}=(\varepsilon+P)u^\mu u^\nu+Pg^{\mu\nu},
\qquad \nabla_\mu T^{\mu\nu}=0.
\]
Nel sistema localmente comovente:
\[
T^{\hat\mu\hat\nu}=\operatorname{diag}(\varepsilon,P,P,P).
\]
$\varepsilon$ include energia a riposo, cinetica e di interazione: è la densità di energia misurata localmente. Se negli appunti del capitolo 2 si usa la segnatura $(+,-,-,-)$, l'espressione covariante diventa $(\varepsilon+P)u^\mu u^\nu-Pg^{\mu\nu}$.
\end{frame}

% Frame 6
\begin{frame}[t]{Metrica sferica statica e massa gravitazionale}
\small
In coordinate areali, con $c=1$, scriviamo
\[
ds^2=-e^{2\Phi(r)}dt^2+
\left(1-\frac{2Gm(r)}r\right)^{-1}dr^2+r^2d\Omega^2.
\]
La funzione $m(r)$ è definita dalla componente radiale della metrica; le equazioni di Einstein danno
\[
\frac{dm}{dr}=4\pi r^2\varepsilon(r).
\]
Per una superficie a $r=R$, $M=m(R)$ è la massa gravitazionale misurata dall'esterno. Il volume proprio di un guscio è
\[
dV_{\rm prop}=4\pi r^2\frac{dr}{\sqrt{1-2Gm(r)/r}},
\]
che \emph{non} è la misura impiegata nell'integrale di $M$.
\end{frame}

% Frame 7
\begin{frame}[t]{Dalla conservazione di $T^{\mu\nu}$ all'equilibrio radiale}
\small
La componente $\nu=r$ dell'equazione $\nabla_\mu T^{\mu\nu}=0$ per un fluido statico è
\[
\frac{dP}{dr}=-(\varepsilon+P)\,\frac{d\Phi}{dr}.
\]
Le equazioni di Einstein per la metrica precedente forniscono
\[
\Phi'(r)=\frac{G[m(r)+4\pi r^3P(r)]}{r[r-2Gm(r)]}.
\]
Combinando:
\[
\frac{dP}{dr}=-\frac{G(\varepsilon+P)[m+4\pi r^3P]}{r(r-2Gm)}.
\]
Questa è l'equazione TOV (Tolman--Oppenheimer--Volkoff).
\end{frame}

% Frame 8
\begin{frame}[t]{Le tre correzioni relativistiche nella TOV}
\small
Riorganizzando la TOV in modo da confrontarla con Newton:
\[
\frac{dP}{dr}=-\frac{Gm\varepsilon}{r^2}
\underbrace{\left(1+\frac{P}{\varepsilon}\right)}_{\text{inerzia entalpica}}
\underbrace{\left(1+\frac{4\pi r^3P}{m}\right)}_{\text{pressione gravita}}
\underbrace{\left(1-\frac{2Gm}{r}\right)^{-1}}_{\text{curvatura}}.
\]
Qui $\varepsilon\simeq\rho$ in unità $c=1$ nel limite newtoniano.
\medskip
Per $P\ll\varepsilon$, $2Gm/r\ll1$ e $4\pi r^3P\ll m$, recuperiamo
$P'\simeq-Gm\rho/r^2$.
\end{frame}

% Frame 9
\begin{frame}[t]{Impostare e integrare una sequenza TOV}
\small
Data una EOS barotropica $\varepsilon=\varepsilon(P)$, fissiamo una pressione centrale $P_c$ e integriamo
\[
\begin{cases}
P'(r)=-G(\varepsilon+P)(m+4\pi r^3P)/[r(r-2Gm)],\\[2pt]
m'(r)=4\pi r^2\varepsilon(P(r)).
\end{cases}
\]
Le condizioni iniziali regolari sono $P(0)=P_c$, $m(0)=0$. Si termina al primo $R$ con $P(R)=0$, ottenendo $M=m(R)$.
\medskip
Ripetendo per diversi $P_c$ si costruiscono $M(P_c)$, $R(P_c)$ e quindi $M(R)$. Un'equazione del primo ordine per il numero barionico serve se si vuole anche $M_B$.
\end{frame}

% Frame 10
\begin{frame}[t]{Sviluppo regolare al centro per iniziare il calcolo numerico}
\small
La TOV non va valutata direttamente con $r=0$ nella forma con denominatore $r(r-2Gm)$. Vicino al centro:
\[
m(r)=\frac{4\pi}{3}\varepsilon_c r^3+O(r^5).
\]
Sostituendo nella TOV con $P(r)=P_c+O(r^2)$:
\[
P'(r)=-\frac{4\pi G}{3}(\varepsilon_c+P_c)(\varepsilon_c+3P_c)r+O(r^3).
\]
Integrando:
\[
P(r)=P_c-\frac{2\pi G}{3}(\varepsilon_c+P_c)(\varepsilon_c+3P_c)r^2+O(r^4).
\]
Si inizializza dunque a $r=r_{\rm start}>0$ con queste espressioni.
\end{frame}

% Frame 11
\begin{frame}[t]{Il salto di energia in una transizione di Maxwell}
\small
Nelle TOV di una stella regolare con $P,\varepsilon>0$, $P'(r)<0$ per $r>0$.
\medskip
Se l'EOS presenta una transizione del primo ordine alla pressione $P_t$, le due densità
\[
\varepsilon_-(P_t)\neq\varepsilon_+(P_t)
\]
sono ammesse, ma la pressione resta continua. Una regione radiale di spessore finito con $P$ perfettamente costante non può sostenere $P'<0$.
\medskip
La costruzione di Maxwell produce così idealmente un'interfaccia a $r=r_t$, non uno strato esteso. È il collegamento con il §9.2.
\end{frame}

\subsection{Neutroni non interagenti}

% Frame 12
\begin{frame}[t]{TOV per neutroni liberi: modello storico}
\small
Il modello di Oppenheimer--Volkoff usa un gas di neutroni non interagenti a $T=0$, la cui EOS è ricavata da
\[
n=\frac{k_F^3}{3\pi^2},\qquad
\varepsilon=\frac{1}{\pi^2}\int_0^{k_F}k^2\sqrt{k^2+m_n^2}\,dk,
\]
\[
P=\frac{1}{3\pi^2}\int_0^{k_F}\frac{k^4}{\sqrt{k^2+m_n^2}}\,dk.
\]
Il massimo riportato dal libro per questo modello è
\[
M_{\max}\simeq0.71M_\odot,\qquad R_{\max}\simeq9.1\ \mathrm{km}.
\]
È una EOS volutamente \emph{irrealistica}: trascura soprattutto le interazioni nucleari che modificano la pressione.
\end{frame}

% Frame 13
\begin{frame}[t]{Figura 4.1: massa e raggio contro densità centrale}
\small
\begin{figure}[!ht]\includegraphics[trim={0cm 0cm 0 0},clip,keepaspectratio,width=0.95\textwidth]{Schaffner/M-R_Pc}\label{fig:Schaffner_M-R_Pc}\end{figure}
La curva della massa raggiunge il primo massimo prima del primo minimo del raggio. La didascalia del libro parla di pressione centrale, ma l'ascissa tracciata è $\varepsilon_c/\varepsilon_0$.
\end{frame}

% Frame 14
\begin{frame}[t]{Figura 4.2: sequenza massa--raggio del gas di neutroni}
\small
\begin{columns}[T,onlytextwidth]
\begin{column}{.62\textwidth}
    \begin{figure}[!ht]\includegraphics[trim={0cm 0cm 0 0},clip,keepaspectratio,width=0.95\textwidth]{Schaffner/MR-freeneutrons}\label{fig:Schaffner_MR-freeneutrons}\end{figure}
\end{column}
\begin{column}{.36\textwidth}
\small
La sequenza parte da masse piccole e raggi grandi; la curva continua è stabile fino al primo massimo di $M$.
\medskip
Il ramo tratteggiato è instabile nel modello del gas libero e tende a una spirale per densità centrali molto alte.
\medskip
La linea inclinata delimita $R=2GM$.
\end{column}
\end{columns}
\end{frame}

% Frame 15
\begin{frame}[t]{Stabilità: che cosa dice il massimo di massa?}
\small
Per una sequenza regolare con EOS fissata, il primo massimo della massa, al crescere di $P_c$, spesso segna la perdita della stabilità del modo radiale fondamentale:
\[
\left.\frac{dM}{dP_c}\right|_{\rm massimo}=0.
\]
Il criterio dei punti di inversione richiede ipotesi su famiglia di equilibri e modi radiali: \emph{un massimo isolato nel diagramma $M(R)$ non basta}. I modelli con transizioni di fase possono mostrare una stabilizzazione successiva (terza famiglia, §9.3).
\medskip
Il gas libero esibisce il comportamento più semplice: ramo stabile prima del massimo e instabile dopo di esso.
\end{frame}

\section{Masse massime}

\subsection{Argomento di Landau}

% Frame 16
\begin{frame}[t]{Argomento di Landau: impostazione e approssimazioni}
\small
Consideriamo $N$ fermioni di massa $m_f$ in una sfera di raggio $R$ e massa totale $M\simeq Nm_f$.
\[
\bar n=\frac{N}{(4\pi/3)R^3},\qquad
k_F=\left(\frac{6\pi^2\bar n}{g}\right)^{1/3}
=\left(\frac{9\pi}{2g}\right)^{1/3}\frac{N^{1/3}}R.
\]
Per neutroni, $g=2$; $k_F$ è il momento caratteristico determinato dal principio di Pauli.
\medskip
Si approssima l'energia gravitazionale per particella come $E_G\sim-GMm_f/R$ e quella cinetica mediante il momento di Fermi. I coefficienti di ordine uno dipendono dal profilo di densità e dalla media sul mare di Fermi.
\end{frame}

% Frame 17
\begin{frame}[t]{Landau: confronto fra degenerazione e gravità}
\small
Nel regime ultrarelativistico del fermione, $E_{\rm kin}\sim k_F$; dunque
\[
E(R)\sim\frac1R\left[
\left(\frac{9\pi}{2g}\right)^{1/3}N^{1/3}
-Gm_f^2N\right].
\]
Entrambi i termini scalano come $R^{-1}$. La possibilità di stabilizzare la sfera cambia al variare di $N$:
\[
Gm_f^2N\sim\left(\frac{9\pi}{2g}\right)^{1/3}N^{1/3}.
\]
La soglia è una stima dimensionale: il coefficiente numerico preciso richiede di risolvere un problema di equilibrio (Newton o TOV) con l'EOS completa.
\end{frame}

% Frame 18
\begin{frame}[t]{Landau: numero critico di fermioni}
\small
Poniamo $a_g=(9\pi/2g)^{1/3}$. Dalla condizione marginale:
\[
Gm_f^2N_{\max}=a_gN_{\max}^{1/3},
\quad N_{\max}^{2/3}=\frac{a_g}{Gm_f^2}.
\]
Elevando alla potenza $3/2$ e usando $G=m_{\rm Pl}^{-2}$:
\[
N_{\max}\sim a_g^{3/2}
\left(\frac{m_{\rm Pl}}{m_f}\right)^3
=\left(\frac{9\pi}{2g}\right)^{1/2}
\left(\frac{m_{\rm Pl}}{m_f}\right)^3.
\]
Questo è l'andamento delle eq. (4.24)--(4.25); la dipendenza $m_f^{-3}$ è più robusta del prefattore.
\end{frame}

% Frame 19
\begin{frame}[t]{Landau: massa e raggio caratteristici}
\small
La massa massima stimata si ottiene moltiplicando per $m_f$:
\[
M_{\max}\sim a_g^{3/2}\frac{m_{\rm Pl}^3}{m_f^2}.
\]
Per il raggio, stimiamo la soglia $k_F\sim m_f$:
\[
R_{\rm crit}\sim a_g\frac{N_{\max}^{1/3}}{m_f}
\sim a_g^{3/2}\frac{m_{\rm Pl}}{m_f^2}.
\]
Definiamo le scale di Landau \emph{senza} prefattori:
\[
M_L\equiv\frac{m_{\rm Pl}^3}{m_f^2},\qquad
R_L\equiv\frac{m_{\rm Pl}}{m_f^2}.
\]
\end{frame}

% Frame 20
\begin{frame}[t]{Stima di Landau vs risultato delle TOV}
\small
Ponendo $m_f=m_n$, il libro riporta
\[
M_L\simeq1.848M_\odot,\qquad R_L\simeq2.729\ \mathrm{km}.
\]
La soluzione numerica per neutroni liberi è invece
\[
M_{\max}=0.71M_\odot,\quad R_{\rm crit}=9.1\ \mathrm{km}.
\]
Nelle soluzioni dimensionless complete (si vedrà nel §4.3) compaiono i coefficienti
\[
M_{\max}\simeq0.384M_L,\qquad R_{\rm crit}\simeq3.367R_L.
\]
La stima $\sim$ di Landau fornisce la \emph{scala}, non il valore numerico esatto.
\end{frame}

\subsection{Limite delle nane bianche}

% Frame 21
\begin{frame}[t]{Nane bianche: chi porta massa e chi porta pressione?}
\small
Per nuclei $(A,Z)$ ed elettroni degeneri, la neutralità elettrica impone
\[
n_e=Zn_A,\qquad N_e=ZN_A=\frac ZA N_B.
\]
La massa totale è dominata dai nuclei:
\[
M\simeq A m_N N_A\simeq m_NN_B,
\]
mentre il sostegno contro gravità proviene principalmente dalla pressione degli elettroni.
\medskip
Quando gli elettroni diventano ultrarelativistici, la loro energia cinetica cresce come $N_e^{4/3}/R$ e può bilanciare la gravità solo fino a un numero barionico critico.
\end{frame}

% Frame 22
\begin{frame}[t]{Limite di Chandrasekhar: dipendenza da $Z/A$}
\small
L'energia cinetica totale degli elettroni è stimata come
\[
U_e\sim N_e k_{F,e}\sim\frac{N_e^{4/3}}R
=\frac{(Z/A)^{4/3}N_B^{4/3}}R.
\]
L'energia gravitazionale totale scala come
\[
U_G\sim-\frac{GM^2}R\sim-\frac{Gm_N^2N_B^2}R.
\]
Alla soglia $|U_G|\sim U_e$:
\[
N_{B,\max}\propto\left(\frac ZA\right)^2
\left(\frac{m_{\rm Pl}}{m_N}\right)^3,
\quad M_{\max}\propto\left(\frac ZA\right)^2\frac{m_{\rm Pl}^3}{m_N^2}.
\]
L'assenza di $m_e$ dalla massa massima è una conseguenza del limite ultrarelativistico degli elettroni.
\end{frame}

% Frame 23
\begin{frame}[t]{Chandrasekhar: risultato numerico e scala del raggio}
\small
Il libro confronta la stima di Landau con il risultato di Chandrasekhar
\[
M_{\rm Ch}\simeq1.44M_\odot\left(\frac{2Z}{A}\right)^2.
\]
Il raggio caratteristico si ricava alla soglia $k_{F,e}\sim m_e$:
\[
R_{\rm WD}\sim\frac{m_{\rm Pl}}{m_Nm_e}
\times\left(\frac ZA\right)
\sim10^3\text{--}10^4\ \mathrm{km}.
\]
In particolare, confrontando le scale di raggio
\[
\frac{R_{\rm WD}}{R_{\rm NS}}\sim\frac{m_N}{m_e}\sim2000
\]
(a fattori composizionali e numerici vicino all'unità).
\end{frame}

% Frame 24
\begin{frame}[t]{Limiti di Landau: quanto sono generali?}
\small
L'argomento dipende dalla competizione fra $U_G\propto-R^{-1}$ e $U_{\rm deg}\propto R^{-1}$ nel regime ultrarelativistico.
\medskip
La stima non determina l'EOS reale: interazioni, composizione, temperatura, relatività generale, neutronizzazione e correzioni Coulomb cambiano i valori critici.
\medskip
Per le nane bianche il caso limite di elettroni ultrarelativistici è ben motivato; per le stelle di neutroni il gas libero è troppo semplice. L'argomento prepara piuttosto la legge di scala comune $m_{\rm Pl}^3/m_f^2$.
\end{frame}

\section{Leggi di scala}

\subsection{Adimensionalizzazione delle TOV}

% Frame 25
\begin{frame}[t]{Scala dimensionale delle equazioni TOV}
\small
Sia $\varepsilon_0$ una scala di densità di energia che compare nell'EOS. Introduciamo
\[
\varepsilon=\varepsilon_0\bar\varepsilon,\quad P=\varepsilon_0\bar P,
\quad r=a\bar r,\quad m=b\bar m.
\]
Cerchiamo $a,b$ tali che le TOV non contengano più coefficienti dimensionali:
\[
\frac{dm}{dr}=4\pi r^2\varepsilon,
\qquad
\frac{dP}{dr}=-\frac{G(\varepsilon+P)(m+4\pi r^3P)}{r(r-2Gm)}.
\]
L'esistenza di una scala unica $\varepsilon_0$ è essenziale: parametri dimensionali aggiuntivi ne alterano la semplice invarianza.
\end{frame}

% Frame 26
\begin{frame}[t]{Equazione della massa: prima condizione sulla scala}
\small
Sostituiamo $r=a\bar r$, $m=b\bar m$, $\varepsilon=\varepsilon_0\bar\varepsilon$:
\[
\frac{b}{a}\frac{d\bar m}{d\bar r}
=4\pi a^2\bar r^2\varepsilon_0\bar\varepsilon.
\]
Vogliamo ottenere la forma adimensionale
\[
\frac{d\bar m}{d\bar r}=4\pi\bar r^2\bar\varepsilon.
\]
Segue la condizione
\[
\frac ba=a^2\varepsilon_0
\quad\Longrightarrow\quad b=a^3\varepsilon_0.
\]
Questa relazione codifica dimensionalmente $m\sim\varepsilon_0r^3$.
\end{frame}

% Frame 27
\begin{frame}[t]{TOV: seconda condizione sulla scala}
\small
Nel fattore di Schwarzschild compare
\[
\frac{2Gm(r)}r=\frac{2Gb}a\frac{\bar m}{\bar r}.
\]
Per recuperare la forma senza costanti scegliamo
\[
\frac{Gb}{a}=1\quad\Longrightarrow\quad b=\frac aG.
\]
Combiniamo con $b=a^3\varepsilon_0$:
\[
\frac aG=a^3\varepsilon_0
\quad\Longrightarrow\quad a^2=\frac1{G\varepsilon_0}.
\]
Dunque
\[
a=(G\varepsilon_0)^{-1/2},\qquad
b=\frac{a}{G}=(G^3\varepsilon_0)^{-1/2}.
\]
\end{frame}

% Frame 28
\begin{frame}[t]{Le TOV adimensionali (4.51)--(4.52)}
\small
Con le scelte precedenti si ottiene
\[
\frac{d\bar m}{d\bar r}=4\pi\bar r^2\bar\varepsilon,
\]
\[
\frac{d\bar P}{d\bar r}=
-\frac{(\bar\varepsilon+\bar P)
(\bar m+4\pi\bar r^3\bar P)}
{\bar r(\bar r-2\bar m)}.
\]
Una EOS della forma $\bar P=f(\bar\varepsilon)$ genera una famiglia di soluzioni adimensionali indipendente dal valore numerico di $\varepsilon_0$.
\medskip
Dopo aver integrato, si recuperano $r=a\bar r$, $m=b\bar m$.
\end{frame}

% Frame 29
\begin{frame}[t]{Leggi di scala per massa, raggio e compattezza}
\small
Dalle scale $a,b$ ricaviamo
\[
R=(G\varepsilon_0)^{-1/2}\bar R,
\qquad M=(G^3\varepsilon_0)^{-1/2}\bar M.
\]
Quindi, mantenendo fissa la \emph{forma} dell'EOS adimensionale,
\[
M_{\max}\propto\varepsilon_0^{-1/2},
\qquad R_{\rm crit}\propto\varepsilon_0^{-1/2}.
\]
La compattezza invece è invariante:
\[
\mathcal C\equiv\frac{GM}{R}=\frac{Gb}{a}\frac{\bar M}{\bar R}
=\frac{\bar M}{\bar R}.
\]
Non tutte le EOS con la stessa $\varepsilon_0$ hanno uguale compattezza: dipende anche dalla \emph{forma} $f$.
\end{frame}

\subsection{Gas libero e EOS del vuoto}

% Frame 30
\begin{frame}[t]{Esempio: il gas libero massivo ha scala $m_f^4$}
\small
Per un fermione libero a $T=0$ ($g=2$) poniamo $x=k_F/m_f$. Dopo il cambio $k=m_ft$, gli integrali hanno la forma
\[
\varepsilon(x)=m_f^4\,\bar\varepsilon(x),\qquad
P(x)=m_f^4\,\bar P(x).
\]
L'unica scala è dunque
\[
\varepsilon_0=m_f^4.
\]
Inserendo nella legge generale:
\[
M_{\max}=\bar M_{\max}\frac{m_{\rm Pl}^3}{m_f^2},
\quad R_{\rm crit}=\bar R_{\rm crit}\frac{m_{\rm Pl}}{m_f^2}.
\]
Si ritrovano esattamente le scale di Landau, questa volta dalle TOV e non dall'argomento energetico semplificato.
\end{frame}

% Frame 31
\begin{frame}[t]{Gas libero: coefficienti e compattezza}
\small
Dall'integrazione delle TOV per il gas libero con $g=2$, i coefficienti riportati nel libro sono
\[
\bar M_{\max}=0.384,\qquad
\bar R_{\rm crit}=3.367.
\]
Pertanto
\[
M_{\max}=0.384\frac{m_{\rm Pl}^3}{m_f^2},
\quad R_{\rm crit}=3.367\frac{m_{\rm Pl}}{m_f^2}.
\]
La compattezza alla massa massima è
\[
\mathcal C_{\max}=\frac{\bar M_{\max}}{\bar R_{\rm crit}}
=\frac{0.384}{3.367}\simeq0.114.
\]
Il valore è indipendente dalla massa del fermione, purché l'EOS sia sempre quella dello stesso gas libero.
\end{frame}

% Frame 32
\begin{frame}[t]{Perché $P=\varepsilon/3$ da sola non fissa una massa massima?}
\small
Per un gas \emph{esattamente} massless e non interagente,
\[
P=\frac13\varepsilon
\]
non compare alcuna scala di energia intrinseca $\varepsilon_0$.
\medskip
Non possiamo ottenere una massa o un raggio caratteristici mediante il solo argomento dimensionale. In più, l'EOS lineare priva di una densità di superficie finita non fornisce la stessa famiglia isolata di stelle autolegate del bag model.
\medskip
Per costruire configurazioni di stelle strange autolegate il libro aggiunge una densità di energia del vuoto.
\end{frame}

% Frame 33
\begin{frame}[t]{Contributo di vuoto: pressione negativa}
\small
Separiamo la parte di gas ultrarelativistico da una costante $B$:
\[
P=P_{\rm free}-B,\qquad
\varepsilon=\varepsilon_{\rm free}+B.
\]
Usando $\varepsilon_{\rm free}=3P_{\rm free}$,
\[
P=\frac{\varepsilon-B}{3}-B
=\frac13(\varepsilon-4B).
\]
Scrivendo $\varepsilon_0=4B$:
\[
P=\frac13(\varepsilon-\varepsilon_0).
\]
La superficie a $P=0$ ha $\varepsilon(R)=\varepsilon_0\neq0$: il sistema è \emph{autolegato}, se la fase considerata è davvero stabile rispetto agli stati alternativi.
\end{frame}

% Frame 34
\begin{frame}[t]{MIT bag: come scala la sequenza massa--raggio}
\small
Per $P=(\varepsilon-4B)/3$ mantenendo il parametro adimensionale $s=1/3$ fissato:
\[
M,R\propto (4B)^{-1/2}\propto B^{-1/2}.
\]
Il libro riporta, con $B^{1/4}=145\ \mathrm{MeV}$ come riferimento,
\[
M_{\max}\simeq2.01M_\odot
\left(\frac{145\,\mathrm{MeV}}{B^{1/4}}\right)^2,
\]
\[
R_{\rm crit}\simeq10.9\,\mathrm{km}
\left(\frac{145\,\mathrm{MeV}}{B^{1/4}}\right)^2.
\]
Sono risultati del modello ideale massless del libro; non garantiscono da soli la stabilità assoluta della strange matter.
\end{frame}

\subsection{Limiti di causalità}

% Frame 35
\begin{frame}[t]{Velocità del suono e rigidità di un\,EOS}
\small
A temperatura zero e composizione in equilibrio, la velocità del suono adiabatico si esprime (in unità $c=1$) come
\[
c_s^2=\frac{dP}{d\varepsilon}.
\]
La causalità impone $0\leq c_s^2\leq1$. Una scelta \emph{massimamente rigida} in questo senso è
\[
\frac{dP}{d\varepsilon}=1,
\]
che, integrata senza costante, dà $P=\varepsilon$ (EOS di Zel'dovich).
\medskip
Questa costruzione è un'estremizzazione, non la previsione che la materia reale raggiunga sempre $c_s=c$.
\end{frame}

% Frame 36
\begin{frame}[t]{EOS di Zel'dovich e derivazione termodinamica}
\small
Consideriamo un contributo energetico dominante quadratico nella densità numerica:
\[
\varepsilon(n)=\kappa n^2.
\]
Al freddo, $P=n^2d(\varepsilon/n)/dn$ e pertanto
\[
P=n^2\frac d{dn}(\kappa n)=\kappa n^2=\varepsilon.
\]
Segue
\[
c_s^2=\frac{dP/dn}{d\varepsilon/dn}=1.
\]
Il risultato vale se il termine $\kappa n^2$ \emph{domina} l'energia; aggiungendo il termine $mn$ la velocità del suono è inferiore a uno a densità finite.
\end{frame}

% Frame 37
\begin{frame}[t]{Proseguire una EOS nota con rigidità causale massima}
\small
Supponiamo di conoscere la EOS fino a $(\varepsilon_f,P_f)$.
Imponiamo, per $\varepsilon>\varepsilon_f$,
\[
\frac{dP}{d\varepsilon}=1.
\]
L'integrazione con $P(\varepsilon_f)=P_f$ dà
\[
P(\varepsilon)=P_f+(\varepsilon-\varepsilon_f),
\qquad\varepsilon\geq\varepsilon_f.
\]
Questa è la prosecuzione più rigida secondo la condizione $c_s^2\leq1$.
\medskip
La parte nucleonica inferiore rimane fissata: il risultato è un \emph{limite superiore} alla massa massima all'interno delle ipotesi sulla regione nota e sulla classe di EOS.
\end{frame}

% Frame 38
\begin{frame}[t]{Rhoades--Ruffini: dipendenza dalla densità fiduciale}
\small
Se $P_f\ll\varepsilon_f$, la scala dominante della prosecuzione causale è $\varepsilon_f$. Allora
\[
M_{\max}\propto\varepsilon_f^{-1/2}.
\]
Il libro cita la stima
\[
M_{\max}\simeq4.2M_\odot
\left(\frac{\varepsilon_{\rm nm}}{\varepsilon_f}\right)^{1/2},
\qquad\varepsilon_{\rm nm}\simeq140\,\mathrm{MeV\,fm^{-3}}.
\]
Per $\varepsilon_f=2\varepsilon_{\rm nm}$ si ottiene $4.2/\sqrt 2\simeq3.0M_\odot$.
\medskip
\emph{Non è un limite causale completamente indipendente dal contesto}: dipende da fino a quale densità l'EOS è considerata nota e da come è raccordata.
\end{frame}

\subsection{EOS autolegate lineari}

% Frame 39
\begin{frame}[t]{EOS lineare autolegata: famiglia a velocità costante}
\small
Una generalizzazione utile del bag model e dell'EOS causale è
\[
P=s\,(\varepsilon-\varepsilon_0),\qquad 0<s\leq1.
\]
La pendenza è costante:
\[
c_s^2=\frac{dP}{d\varepsilon}=s.
\]
$\varepsilon_0$ è la densità di energia alla superficie, dove $P=0$.
\medskip
Per $s=1/3$ e $\varepsilon_0=4B$ si ottiene la EOS del bag ideale; per $s=1$ si ha la massima rigidità causale sopra una densità di superficie finita.
\end{frame}

% Frame 40
\begin{frame}[t]{Scalare $\varepsilon_0$ non equivale a cambiare $s$}
\small
Per una EOS lineare autolegata,
\[
M_{\max}=\mathcal M(s)(G^3\varepsilon_0)^{-1/2},
\qquad R_{\rm crit}=\mathcal R(s)(G\varepsilon_0)^{-1/2}.
\]
I coefficienti adimensionali $\mathcal M(s),\mathcal R(s)$ devono essere ottenuti integrando le TOV e \emph{dipendono da $s$}.
\[
\frac{GM_{\max}}{R_{\rm crit}}
=\frac{\mathcal M(s)}{\mathcal R(s)}.
\]
Variare $\varepsilon_0$ riscalando massa e raggio non modifica la forma adimensionale della curva; variare $s$ sì.
\end{frame}

% Frame 41
\begin{frame}[t]{Esempi quantitativi: $s=1/3,\,2/3,\,1$}
\small
Il libro (§4.3.5) usa $\varepsilon_{\rm nm}=140\,\mathrm{MeV\,fm^{-3}}$ come densità di riferimento:
\[
\begin{array}{c|ccc}
s &1/3&2/3&1\\\hline
M_{\max}/(M_\odot\sqrt{\varepsilon_{\rm nm}/\varepsilon_0})
&2.57&3.64&4.23\\
R_{\rm crit}/(\mathrm{km}\sqrt{\varepsilon_{\rm nm}/\varepsilon_0})
&14.0&16.4&17.6
\end{array}
\]
All'aumentare di $s$, la massa massima cresce; i raggi non seguono però la stessa legge in funzione di $s$.
\end{frame}

% Frame 42
\begin{frame}[t]{Compattezza: i valori della eq. (4.80)}
\small
La compattezza $\mathcal C=GM/R$ non dipende dalla scala $\varepsilon_0$ per una forma fissata di EOS.
\[
\begin{array}{c|ccc}
s&1/3&2/3&1\\\hline
\mathcal C_{\max}&0.271&0.328&0.354
\end{array}
\]
Per confronto, nel modello a neutroni liberi si ha $\mathcal C_{\max}\simeq0.11$.
\medskip
Questi valori sono quelli della famiglia autolegata lineare. Il valore $0.354$ viene anche utilizzato dal libro come riferimento per limiti causali più generali, sotto ipotesi aggiuntive.
\end{frame}

% Frame 43
\begin{frame}[t]{Redshift gravitazionale dalla superficie}
\small
All'esterno della stella vale la metrica di Schwarzschild. Per un fotone emesso dalla superficie e misurato all'infinito,
\[
1+z=\frac{1}{\sqrt{1-2GM/R}}=\frac{1}{\sqrt{1-2\mathcal C}}.
\]
Ai tre valori di compattezza massima della famiglia lineare il libro associa
\[
\begin{array}{c|ccc}
s&1/3&2/3&1\\\hline
z_{\max}&0.478&0.705&0.851
\end{array}
\]
Per $\mathcal C\ll1$, espandendo la radice si recupera $z\simeq\mathcal C=GM/R$, il limite debole discusso nel capitolo 2.
\end{frame}

% Frame 44
\begin{frame}[t]{Densità centrale massima della famiglia autolegata}
\small
La densità centrale alla configurazione di massa massima, divisa per la densità di superficie $\varepsilon_0$, è adimensionale:
\[
\begin{array}{c|ccc}
s&1/3&2/3&1\\\hline
\varepsilon_{c,\max}/\varepsilon_0&4.81&3.54&3.03
\end{array}
\]
Per valori diversi di $\varepsilon_0$, le densità si riscalano linearmente, mentre masse e raggi scalano come $\varepsilon_0^{-1/2}$.
\medskip
Non si deve confondere la densità \emph{centrale} con la densità alla \emph{superficie} $\varepsilon_0$.
\end{frame}

% Frame 45
\begin{frame}[t]{Relazione di Lattimer--Prakash: derivazione della scala}
\small
Per una famiglia a forma dimensionless fissata:
\[
M_{\max}=A\varepsilon_0^{-1/2},
\qquad\varepsilon_{c,\max}=B\varepsilon_0.
\]
Eliminando $\varepsilon_0$, si ottiene
\[
\varepsilon_{c,\max}M_{\max}^2=BA^2=\text{costante}.
\]
Quindi
\[
\varepsilon_{c,\max}\propto M_{\max}^{-2}.
\]
Il libro cita, per la costruzione causale limite:
\[
\varepsilon_{\rm ultimate}\simeq13.5\varepsilon_{\rm nm}
\left(\frac{2M_\odot}{M_{\max}}\right)^2.
\]
La relazione esprime il massimo caratteristico all'interno delle ipotesi della famiglia, non una densità universale di qualunque materia.
\end{frame}

% Frame 46
\begin{frame}[t]{Esempio numerico della relazione massa--densità}
\small
Con $\varepsilon_{\rm nm}\simeq140\,\mathrm{MeV\,fm^{-3}}$,
\[
M_{\max}=2M_\odot\quad\Longrightarrow\quad
\varepsilon_{\rm ultimate}\simeq13.5\varepsilon_{\rm nm}.
\]
Per una massa massima ipotetica di $3M_\odot$:
\[
\varepsilon_{\rm ultimate}
\simeq13.5\left(\frac23\right)^2\varepsilon_{\rm nm}
=6\varepsilon_{\rm nm}.
\]
Una massa massima maggiore richiede, in questa costruzione, una scala energetica/densità minore.
\medskip
Riferimento concettuale: la massa stellare osservata vincola indirettamente la regione ad alte densità dell'EOS.
\end{frame}

% Frame 47
\begin{frame}[t]{Riepilogo delle quattro EOS del §4.3}
\small
\begin{tabular}{@{}lll@{}}
\textbf{EOS} & \textbf{Scala} & \textbf{Interpretazione}\\\hline
$P=P_{\rm free}(m_f)$ & $m_f^4$ & Fermi gas massivo\\
$P=(\varepsilon-4B)/3$ & $B$ & Bag, autolegata\\
$P=P_f+\varepsilon-\varepsilon_f$ & $\varepsilon_f$ & Prosecuzione causale\\
$P=s(\varepsilon-\varepsilon_0)$ & $\varepsilon_0,s$ & Autolegata generica
\end{tabular}
\bigskip
La legge $M,R\propto\varepsilon_0^{-1/2}$ è valida quando tutte le caratteristiche dimensionali della EOS scalano coerentemente con l'unica densità $\varepsilon_0$.
\medskip
In una EOS composta di crosta, core, transizioni e soglie di particelle, la riscalatura può non essere esatta.
\end{frame}

\section{Fermioni interagenti}

\subsection{Interazione quadratica}

% Frame 48
\begin{frame}[t]{Interazioni fra fermioni: ansatz del libro}
\small
Aggiungiamo all'energia di un gas di fermioni un termine efficace
\[
\varepsilon_{\rm int}(n)=c\,n^\gamma.
\]
Per la termodinamica a $T=0$,
\[
P=n^2\frac{d}{dn}\left(\frac\varepsilon n\right).
\]
Il contributo di interazione alla pressione è
\[
P_{\rm int}=n^2\frac{d}{dn}(c n^{\gamma-1})
=(\gamma-1)c n^\gamma=(\gamma-1)\varepsilon_{\rm int}.
\]
Quando \emph{domina} tale contributo, $c_s^2\to\gamma-1$; l'ansatz sarebbe acausale a densità elevate per $\gamma>2$.
\end{frame}

% Frame 49
\begin{frame}[t]{Caso $\gamma=2$: energia e pressione}
\small
Per un termine repulsivo quadratico prendiamo $c=1/m_I^2>0$:
\[
\varepsilon_{\rm int}=\frac{n^2}{m_I^2},\qquad
P_{\rm int}=\frac{n^2}{m_I^2}.
\]
Nella regione non relativistica per i fermioni,
\[
\varepsilon(n)\simeq m_fn+\varepsilon_{\rm kin}^{\rm NR}(n)+\frac{n^2}{m_I^2},
\]
\[
P(n)\simeq P_{\rm kin}^{\rm NR}(n)+\frac{n^2}{m_I^2},
\qquad P_{\rm kin}^{\rm NR}\propto\frac{n^{5/3}}{m_f}.
\]
La costante $m_I$ parametrizza una scala di interazione \emph{effettiva}, non necessariamente la massa di una particella reale.
\end{frame}

% Frame 50
\begin{frame}[t]{Tre regimi della EOS con interazione repulsiva}
\small
Per $n$ crescente, l'ansatz può attraversare tre comportamenti:
\begin{enumerate}
\item Bassa densità: $\varepsilon\simeq m_fn$ e $P\propto n^{5/3}$, dunque $P\propto\varepsilon^{5/3}$.
\item Regime intermedio: $\varepsilon\simeq m_fn$, ma $P\simeq n^2/m_I^2$, dunque $P\propto\varepsilon^2$.
\item Interazione dominante anche in $\varepsilon$: $P\simeq\varepsilon\simeq n^2/m_I^2$, con $c_s^2\to1$.
\end{enumerate}
La rigidezza massima a densità molto grandi emerge come limite dell'ansatz quadratico, non è garantita per la QCD reale.
\end{frame}

% Frame 51
\begin{frame}[t]{Interpretazione RMF del termine quadratico}
\small
Nella teoria relativistica di campo medio (RMF) con campo vettoriale $\omega_\mu$, l'equazione del campo statico è
\[
m_\omega^2\omega_0=g_\omega n_B.
\]
L'energia vettoriale repulsiva vale schematicamente
\[
\varepsilon_\omega=\frac12m_\omega^2\omega_0^2
=\frac{g_\omega^2}{2m_\omega^2}n_B^2,
\]
e produce un contributo alla pressione $P_\omega=\varepsilon_\omega$ nel modello lineare.
\medskip
È il collegamento con i nostri appunti RMF/Walecka, ma non riproduce da solo la EOS nucleonica completa (campi scalari, simmetria, composizione, ecc.).
\end{frame}

% Frame 52
\begin{frame}[t]{Parametro adimensionale di interazione $y$}
\small
Il libro introduce
\[
y\equiv\frac{m_f}{m_I},\qquad c=\frac1{m_I^2}.
\]
In unità del fermione ($n=m_f^3\bar n$), l'interazione contribuisce
\[
\frac{\varepsilon_{\rm int}}{m_f^4}=
\frac{n^2}{m_I^2m_f^4}=y^2\bar n^2.
\]
Quindi
\[
y\ll1:\ \text{interazione debole},\qquad
y\gg1:\ \text{interazione forte}
\]
\emph{nel senso specifico di questo modello}. Non si può identificare direttamente $y$ con $\alpha_s$ della QCD.
\end{frame}

% Frame 53
\begin{frame}[t]{Densità di crossover e scala di massa interagente}
\small
Per $y\gg1$, la transizione fra energia di riposo e interazione si stima imponendo
\[
m_fn_c\sim\frac{n_c^2}{m_I^2}
\quad\Longrightarrow\quad n_c\sim m_fm_I^2
=\frac{m_f^3}{y^2}.
\]
Allora la scala energetica della EOS è
\[
\varepsilon_c\sim m_fn_c\sim\frac{m_f^4}{y^2}.
\]
Dalla riscalatura delle TOV:
\[
M_{\max}\sim\frac{m_{\rm Pl}^3}{\sqrt{\varepsilon_c}}
\sim y\frac{m_{\rm Pl}^3}{m_f^2}=yM_L,
\]
\[
R_{\rm crit}\sim yR_L.
\]
Le dipendenze lineari richiedono il regime di interazioni sufficientemente forti.
\end{frame}

% Frame 54
\begin{frame}[t]{Fit numerici del libro per i fermioni interagenti}
\small
Per la particolare famiglia di EOS del §4.4 il testo riporta i fit
\[
M_{\max}(y)\simeq(0.384+0.165y)M_L,
\]
\[
R_{\rm crit}(y)\simeq(3.367+0.797y)R_L.
\]
A $y=0$ recuperiamo i coefficienti TOV del gas libero.
\medskip
Per $m_f=m_n$ e $y\sim10$ la stima dà una massa massima dell'ordine di $3.8M_\odot$, ma \emph{è un'illustrazione del modello parametrico}, non una previsione quantitativa affidabile per stelle di neutroni reali.
\end{frame}

% Frame 55
\begin{frame}[t]{Figura 4.3: curve massa--raggio con interazioni}
\small
\begin{columns}[T,onlytextwidth]
\begin{column}{.62\textwidth}
    \begin{figure}[!ht]\includegraphics[trim={0cm 0cm 0 0},clip,keepaspectratio,width=0.95\textwidth]{Schaffner/MR-interactingfermions_y}\label{fig:Schaffner_MR-interactingfermions_y}\end{figure}
\end{column}
\begin{column}{.36\textwidth}
\small
Con $y\ll1$ le curve si avvicinano al comportamento del gas libero.
\medskip
Con $y\gg1$ crescono massa massima e raggio caratteristico; compare un tratto quasi a raggio costante.
\medskip
Gli assi rappresentano grandezze adimensionali $M'=M/M_L$ e $R'=R/R_L$ in scala logaritmica.
\end{column}
\end{columns}
\end{frame}

% Frame 56
\begin{frame}[t]{Compattezza del modello interagente e limiti}
\small
Dal rapporto dei fit numerici,
\[
\mathcal C_{\max}(y)
=\frac{0.384+0.165y}{3.367+0.797y}.
\]
Per $y\to0$:
\[
\mathcal C_{\max}\to0.384/3.367\simeq0.114.
\]
Per $y\to\infty$:
\[
\mathcal C_{\max}\to0.165/0.797\simeq0.207.
\]
Il modello illustra come la repulsione possa aumentare la compattezza massima rispetto al caso libero. Non è un limite causale universale: dipende dalla forma scelta per l'interazione.
\end{frame}

% Frame 57
\begin{frame}[t]{Interazioni e fisica nucleare reale}
\small
Il libro confronta scale efficaci: $m_I\sim f_\pi\simeq92\,\mathrm{MeV}$ porta formalmente a $y\sim m_n/f_\pi\sim10$.
\medskip
Tuttavia gli accoppiamenti, i canali attrattivi e repulsivi, la densità di simmetria, le correzioni a molti corpi e le nuove specie rendono la EOS reale più articolata.
\medskip
L'ansatz $+cn^2$ e il legame con $\omega_0$ hanno soprattutto valore \emph{esplicativo}. Le EOS di stelle di neutroni vanno ottenute da modelli microscopici verificati e confrontate con $M_{\max}\gtrsim2M_\odot$ e altri vincoli.
\end{frame}

\section{Stelle bosoniche}

\subsection{Bosoni liberi}

% Frame 58
\begin{frame}[t]{Stelle bosoniche: perché non sono gas di neutroni}
\small
Nel §4.5 il libro considera \emph{stelle composte da particelle bosoniche ipoteticamente stabili}, non la superfluidità neutronica BCS studiata negli appunti del capitolo 3.
\medskip
La superfluidità dei neutroni modifica eccitazioni e dinamica del fluido, ma non trasforma automaticamente una stella di neutroni in una ``boson star'' fondamentale.
\medskip
Una boson star realistica è descritta da un campo scalare autogravitante e dalle equazioni di Einstein--Klein--Gordon; un modello di fluido barotropico può approssimarne alcuni regimi, soprattutto con autointerazioni repulsive dominanti.
\end{frame}

% Frame 59
\begin{frame}[t]{Bosoni liberi: principio di indeterminazione}
\small
A $T=0$ bosoni non interagenti possono occupare lo stato a momento minimo; la pressione di degenerazione di Pauli non esiste.
\medskip
Per un oggetto quantistico confinato entro un raggio $R$, però, il gradiente spaziale del campo implica un momento caratteristico
\[
\Delta p\sim\frac1R.
\]
Stimiamo l'energia cinetica per particella nei due limiti:
\[
E_{\rm kin}^{\rm UR}\sim\frac1R,
\qquad E_{\rm kin}^{\rm NR}\sim\frac1{m_bR^2}.
\]
Sono argomenti di scala legati alla dimensione del sistema, non una pressione termica classica.
\end{frame}

% Frame 60
\begin{frame}[t]{Bosoni liberi: bilancio energetico di Landau}
\small
Per $N$ bosoni di massa $m_b$ e $M\simeq Nm_b$,
\[
E(R)\sim-\frac{GMm_b}R+\frac1R
=\frac{1-GNm_b^2}{R}
\]
nel limite UR. Alla soglia marginale
\[
GN_{\max}m_b^2\sim1,
\quad N_{\max}\sim\frac{m_{\rm Pl}^2}{m_b^2}.
\]
Quindi
\[
M_{\max}^{\rm free}\sim\frac{m_{\rm Pl}^2}{m_b},
\qquad R_{\rm crit}^{\rm free}\sim\frac1{m_b}.
\]
La soluzione di campo completa fornisce coefficienti numerici e profili dipendenti dall'ansatz del campo.
\end{frame}

% Frame 61
\begin{frame}[t]{Perché il limite dei bosoni liberi differisce da Landau?}
\small
Confrontiamo le dipendenze alla stessa massa microscopica $m$:
\[
M_{\max}^{\rm fermioni}\sim\frac{m_{\rm Pl}^3}{m^2},
\qquad M_{\max}^{\rm bosoni\ liberi}\sim\frac{m_{\rm Pl}^2}{m}.
\]
Il rapporto è
\[
\frac{M_{\max}^{\rm fermioni}}{M_{\max}^{\rm bosoni\ liberi}}
\sim\frac{m_{\rm Pl}}m\gg1
\]
per particelle con $m\ll m_{\rm Pl}$. I fermioni possiedono un enorme sostegno dalla degenerazione; i bosoni liberi dipendono dalla struttura spaziale del campo e dal principio d'indeterminazione.
\end{frame}

\subsection{Autointerazioni e stelle bosoniche}

% Frame 62
\begin{frame}[t]{Bosoni con interazione repulsiva: EOS del libro}
\small
Consideriamo ora il modello
\[
\varepsilon(n)=m_bn+cn^2,\qquad c>0.
\]
Per consistenza termodinamica
\[
P=n^2\frac{d}{dn}\left(m_b+cn\right)=cn^2.
\]
A basse densità $\varepsilon\simeq m_bn$ e quindi
\[
P\simeq\frac c{m_b^2}\varepsilon^2
\]
(approssimazione politropica con indice $2$).
\medskip
Ad alte densità $\varepsilon\simeq cn^2$ e quindi $P\simeq\varepsilon$: comportamento causalmente rigido nell'ansatz.
\end{frame}

% Frame 63
\begin{frame}[t]{Bosoni interagenti: scala critica e massa massima}
\small
Scriviamo $c=1/m_I^2$ e $y=m_b/m_I$. Il punto in cui massa a riposo e interazioni sono comparabili soddisfa
\[
m_bn_c\sim\frac{n_c^2}{m_I^2}
\ \Longrightarrow\ n_c\sim m_bm_I^2
=\frac{m_b^3}{y^2}.
\]
La densità di energia caratteristica è
\[
\varepsilon_c\sim m_bn_c\sim\frac{m_b^4}{y^2}.
\]
Dalla legge di scala TOV:
\[
M_{\max}\sim y\frac{m_{\rm Pl}^3}{m_b^2},
\qquad R_{\rm crit}\sim y\frac{m_{\rm Pl}}{m_b^2}.
\]
Questo regime è differente dal caso libero e richiede interazioni sufficientemente importanti.
\end{frame}

% Frame 64
\begin{frame}[t]{Bosoni interagenti: prefattori riportati dal libro}
\small
Nel modello efficace del §4.5, il testo dà
\[
M_{\max}\simeq0.164\,y\,M_L,
\qquad R_{\rm crit}\simeq0.763\,y\,R_L,
\]
con le scale $M_L=m_{\rm Pl}^3/m_b^2$ e $R_L=m_{\rm Pl}/m_b^2$.
\medskip
La compattezza al massimo vale
\[
\mathcal C_{\max}\simeq\frac{0.164}{0.763}\simeq0.215.
\]
I coefficienti dipendono dal modello di materia bosonica adottato; la soluzione completa richiede di precisare campo, potenziale e assunzioni di isotropia.
\end{frame}

% Frame 65
\begin{frame}[t]{Figura 4.4: massa--raggio per bosoni con autointerazione}
\small
\begin{columns}[T,onlytextwidth]
\begin{column}{.62\textwidth}
    \begin{figure}[!ht]\includegraphics[trim={0cm 0cm 0 0},clip,keepaspectratio,width=0.95\textwidth]{Schaffner/MR-interactingboson_y}\label{fig:Schaffner_MR-interactingboson_y}\end{figure}
\end{column}
\begin{column}{.36\textwidth}
\small
Le curve sono rappresentate in scala logaritmica per vari $y$.
\medskip
In un ampio intervallo di masse il raggio resta quasi costante, come atteso dalla politropica $P\propto\varepsilon^2$.
\medskip
La crescita di $y$ amplifica le scale di massa e raggio.
\end{column}
\end{columns}
\end{frame}

% Frame 66
\begin{frame}[t]{Quando servono le equazioni di Einstein--Klein--Gordon?}
\small
Una boson star di campo scalare complesso $\phi$ è tipicamente descritta mediante l'azione
\[
S=\int d^4x\,\sqrt{-g}\left[
\frac{R}{16\pi G}-g^{\mu\nu}\partial_\mu\phi^*\partial_\nu\phi
-V(|\phi|^2)\right].
\]
Lo stress--energia contiene anche i \emph{gradienti} del campo. Il modello perfetto-fluido $P=P(\varepsilon)$ non li contiene e non è esatto per un condensato bosonico autogravitante in generale.
\medskip
La rappresentazione EOS del libro è soprattutto uno strumento utile per stimare le scale. Le configurazioni di Kaup e i modelli con autointerazioni sono soluzioni di campo più fondamentali.
\end{frame}

% Frame 67
\begin{frame}[t]{Confronto e limiti di applicabilità dei modelli}
\small
\begin{tabular}{@{}lll@{}}
\textbf{Sistema} & \textbf{Sostegno principale} & \textbf{Scala $M_{\max}$}\\\hline
Fermioni liberi & Degenerazione & $m_{\rm Pl}^3/m_f^2$\\
Fermioni interagenti & Degenerazione + repulsione & $\sim yM_L$\\
Bosoni liberi & Gradiente quantistico & $m_{\rm Pl}^2/m_b$\\
Bosoni interagenti & Autointerazione & $\sim y m_{\rm Pl}^3/m_b^2$
\end{tabular}
\bigskip
Tutte sono \emph{stime di scala o modelli idealizzati}, non classificazioni esaustive delle possibili stelle compatte.
\end{frame}

\section{Sintesi e collegamenti}

% Frame 68
\begin{frame}[t]{Riepilogo delle derivazioni fondamentali del capitolo 4}
\small
\begin{enumerate}
\item Equilibrio: $P'=-G(\varepsilon+P)(m+4\pi r^3P)/[r(r-2Gm)]$ e $m'=4\pi r^2\varepsilon$.
\item Landau: $M_L\sim m_{\rm Pl}^3/m_f^2$, $R_L\sim m_{\rm Pl}/m_f^2$.
\item Scala TOV: $M\propto\varepsilon_0^{-1/2}$, $R\propto\varepsilon_0^{-1/2}$ per EOS a scala unica.
\item Causalità: $dP/d\varepsilon\leq1$ e prosecuzione estrema $P=P_f+\varepsilon-\varepsilon_f$.
\item Interazioni: $\varepsilon_{\rm int}=cn^\gamma$, $P_{\rm int}=(\gamma-1)cn^\gamma$.
\item Stelle bosoniche: differenza fra scala libera $m_{\rm Pl}^2/m_b$ e scala con autointerazione.
\end{enumerate}
\end{frame}

% Frame 69
\begin{frame}[t]{Che cosa riprendere dai capitoli successivi?}
\small
Il capitolo 4 fornisce le basi comuni per:
\begin{itemize}
\item §5: Lane--Emden e nane bianche, limite di Chandrasekhar.
\item §7: stelle di neutroni con EOS nucleoniche realistiche, repulsione, simmetria e nuove specie.
\item §8: stelle autolegate di strange matter e MIT bag.
\item §9: transizioni del primo ordine, criterio di Seidov, terza famiglia e bande massa--raggio.
\end{itemize}
Il risultato più generale è che la gravità determina l'equilibrio una volta data $P(\varepsilon)$, ma l'identità microscopica della materia non può essere letta direttamente da una singola misura di $M$ e $R$.
\end{frame}

